\documentclass[10pt,a4paper]{article}
\usepackage{nexusS5}
\setnexusheader{Entrée en Quatrième — Stage de pré-rentrée}{Séance 5/5 — Sinda CHIKHAOUI}
\begin{document}
\nexuscover{SÉANCE 5 --- RÉINVESTIR, CONSOLIDER, PRÉPARER LA RENTRÉE}{Livret de travail}{Sinda CHIKHAOUI}{Transfert et plan de rentrée}{Entrée en Quatrième --- Mathématiques --- 1\,h\,15 de travail, puis 45 minutes d'évaluation dans un document séparé}
\begin{goalbox}
\textbf{Objectif personnel de cette séance.} Stabiliser la soustraction d'un nombre négatif et la réduction d'une expression comportant des constantes négatives, deux points nommés dans le plan de remédiation.
\par\medskip
\textbf{Ce livret couvre uniquement les 75 premières minutes.} L'évaluation finale est remise ensuite, dans un document distinct.
\end{goalbox}
\begin{methodbox}
\textbf{Le fil des cinq séances}
\begin{itemize}[leftmargin=13mm,labelsep=3mm,itemsep=1.5pt]
\item[\textbf{\color{Navy}S1}] Calculer avec du sens (relatifs, fractions) ;
\item[\textbf{\color{Navy}S2}] Mesurer une surface ou un contour (aires et périmètres) ;
\item[\textbf{\color{Navy}S3}] Donner du sens aux lettres (calcul littéral) ;
\item[\textbf{\color{Navy}S4}] Voir, raisonner, démontrer (géométrie) ;
\item[\textbf{\color{Navy}S5}] aujourd'hui : réactiver, consolider deux axes personnels, faire le pont vers la Quatrième, puis mesurer.
\end{itemize}
\end{methodbox}
\sectiontitle{Feuille de route des 75 minutes}
\rowcolors{2}{SoftBlue}{white}
\par\noindent\begin{tabularx}{\linewidth}{>{\bfseries}p{20mm} X X}
\rowcolor{Navy}\color{white}Temps & \color{white}Travail & \color{white}Preuve attendue \\
0--10 min & Remettre en activité sur les cinq domaines du stage et faire apparaître immédiatement les oublis. & Cinq réponses écrites et le contrôle utilisé. \\
10--35 min & Somme, différence et soustraction d'un nombre négatif & Trois réussites autonomes sur l'axe prioritaire. \\
35--55 min & Réduire une expression en respectant le signe des constantes & Deux réussites autonomes sur le second axe. \\
55--70 min & De la Cinquième à la Quatrième : le signe d'un produit et la première équation & Une production reliant un acquis de Cinquième à un attendu de Quatrième. \\
70--75 min & Cadrage méthodologique, sans aucune information sur le contenu de l'évaluation. & Deux engagements écrits avant l'évaluation. \\
\end{tabularx}
\begin{graybox}\small\textbf{Règle de la séance.} J'écris ma réponse \emph{et} le contrôle qui la rend défendable. Une réponse sans contrôle reste une hypothèse. La ligne « Certitude » se remplit \emph{avant} toute vérification.\end{graybox}
\par\vspace{2mm}
\phasehead{Phase 1 --- Réactivation}{0--10 min}{Remettre en activité sur les cinq domaines du stage et faire apparaître immédiatement les oublis.}
Sept minutes seules, sans carte de rappel, puis trois minutes de mise en commun. Écrire la réponse \emph{et} le contrôle utilisé.
\begin{enumerate}
\item Calculer $(-9) + 4 - (-6)$.
\shortlines{2}
\item Calculer $\dfrac{7}{10} - \dfrac{2}{5}$ en transformant entièrement la seconde fraction.
\shortlines{2}
\item Réduire $6x - 5 + 2x + 9$, puis contrôler pour $x = 1$.
\shortlines{2}
\item Deux angles d'un triangle mesurent $43^\circ$ et $68^\circ$. Calculer le troisième en citant la propriété.
\shortlines{2}
\item Un rectangle mesure 11 cm sur 6 cm. Donner son aire et son périmètre, unités comprises.
\shortlines{2}
\end{enumerate}
\certline
\aideline
\needspace{62mm}\par\vspace{3mm}
\phasehead{Phase 2 --- Consolidation, axe prioritaire}{10--35 min}{Somme, différence et soustraction d'un nombre négatif}
\begin{methodbox}
\textbf{Une seule règle.} Soustraire un nombre, c'est ajouter son opposé : $a - (-b) = a + b$.
Sur la droite graduée, soustraire un nombre négatif revient à se déplacer vers la \textbf{droite}.
\textbf{Contrôle :} un résultat doit toujours pouvoir être situé approximativement sur la droite graduée.
\end{methodbox}
\begin{graybox}
\textbf{Exemple guidé.} Calculer $6 - (-9) + (-4)$.
\begin{enumerate}
\item Réécrire chaque soustraction en addition : $6 + 9 + (-4)$.
\item Calculer de gauche à droite : $15 + (-4) = 11$.
\item Contrôler : on part de $6$, on avance de $9$ (donc au-delà de $10$), puis on recule de $4$.
\end{enumerate}
\end{graybox}
\subsectiontitle{À toi}
\begin{enumerate}
\needspace{18mm}
\item Calculer $8 - (-3)$.
\worklines{2}
\needspace{18mm}
\item Calculer $-6 + (-5) - (-9)$.
\worklines{2}
\needspace{18mm}
\item Calculer $12 - (-4) - 7$.
\worklines{2}
\needspace{18mm}
\item Un point d'abscisse $-11$ se déplace de $15$ unités vers la droite. Donner son abscisse finale.
\worklines{2}
\needspace{18mm}
\item Compléter : $-7 - (\ldots) = -2$, puis expliquer la démarche.
\worklines{3}
\needspace{18mm}
\item Un élève écrit $9 - (-6) = 3$. Nommer précisément son erreur, donner le résultat exact et proposer un contrôle.
\worklines{4}
\end{enumerate}
\begin{successbox}\textbf{Critère de réussite de cette phase :} quatre calculs exacts consécutifs, dont deux comportant une soustraction d'un nombre négatif, avec un contrôle écrit.\end{successbox}
\certline
\needspace{62mm}\par\vspace{3mm}
\phasehead{Phase 3 --- Consolidation, second axe}{35--55 min}{Réduire une expression en respectant le signe des constantes}
\begin{methodbox}
\textbf{Deux familles, jamais mélangées.} Les termes en $x$ se regroupent entre eux, les constantes entre elles.
Le signe qui précède un nombre lui \textbf{appartient} : dans $3 - 8$, la constante est $-8$.
\textbf{Contrôle :} remplacer $x$ par $0$ ne laisse que les constantes ; remplacer $x$ par $1$ vérifie l'ensemble.
\end{methodbox}
\subsectiontitle{À toi}
\begin{enumerate}
\needspace{18mm}
\item Réduire $8x + 3 - 5x - 11$.
\worklines{2}
\needspace{18mm}
\item Réduire $-4x + 9 + 7x - 2$.
\worklines{2}
\needspace{18mm}
\item Développer puis réduire $3(x + 4) + 2x - 5$.
\worklines{3}
\needspace{18mm}
\item Développer puis réduire $6(x - 2) - x + 9$.
\worklines{3}
\end{enumerate}
\begin{successbox}\textbf{Critère de réussite de cette phase :} trois réductions consécutives exactes avec constantes signées, dont une après développement, chacune contrôlée par substitution.\end{successbox}
\certline
\needspace{62mm}\par\vspace{3mm}
\phasehead{Phase 4 --- Réinvestissement}{55--70 min}{Découverte guidée : construire, à partir d'acquis de Cinquième, deux notions nouvelles de Quatrième — le signe d'un produit de relatifs et l'équation du premier degré.}

\begin{methodbox}
\textbf{Ce qui change en Quatrième.} L'addition et la soustraction des relatifs sont supposées acquises :
elles ont été travaillées en séance 1. La Quatrième y ajoute le \textbf{produit} et le \textbf{quotient} de
relatifs, puis utilise les expressions littérales pour écrire et résoudre des \textbf{équations}.
\par\smallskip
Ces deux notions sont \textbf{nouvelles}. Les blocs ci-dessous ne sont donc pas un réemploi : ce sont deux
\textbf{découvertes guidées}. Chacune part d'un acquis de Cinquième — le sens des signes, l'écriture d'une
expression littérale — pour construire la règle nouvelle pas à pas.
\end{methodbox}
\begin{warningbox}\small
\textbf{Cette phase introduit des notions de l'année à venir.} Ce que tu produis ici ne sera donc
\emph{pas} comparé à ton positionnement de début de stage : il ne s'agit pas de mesurer un progrès,
mais de préparer les premières séances de septembre.
\end{warningbox}


\subsectiontitle{A. Le signe d'un produit — 6 min}
Compléter le tableau, puis écrire la règle en une phrase.
\par\noindent\begin{tabularx}{\linewidth}{>{\bfseries}p{30mm} X >{\bfseries}p{30mm} X}
\toprule
Produit & Résultat & Produit & Résultat \\
\midrule
$3 \times (-4)$ & & $(-3)\times(-4)$ & \\
$(-5)\times 2$ & & $(-5)\times(-2)$ & \\
\bottomrule
\end{tabularx}
\par\vspace{1mm}
Ma règle : \rule{125mm}{0.32pt}
\par\vspace{1mm}
Contrôle demandé : donner un produit de deux nombres dont le résultat est $-12$, puis un autre, différent.
\shortlines{1}

\subsectiontitle{B. Écrire puis résoudre une équation — 9 min}
Une salle est louée $35$ TND, auxquels s'ajoutent $4$ TND par participant.
\begin{enumerate}
\item Écrire l'expression $C(n)$ du coût pour $n$ participants. \rule{60mm}{0.32pt}
\item Calculer $C(12)$ en montrant la substitution.
\shortlines{1}
\item La facture s'élève à $99$ TND. Écrire l'équation correspondante, puis la résoudre en écrivant chaque étape.
\worklines{3}
\item Contrôler la solution trouvée en la remplaçant dans $C(n)$.
\shortlines{1}
\end{enumerate}

\begin{successbox}
\textbf{Preuve attendue :} l'expression est écrite avant tout calcul ; la solution de l'équation est
contrôlée par substitution et non par une simple relecture.
\end{successbox}

\needspace{62mm}\par\vspace{3mm}
\phasehead{Phase 5 --- Avant l'évaluation}{70--75 min}{Cadrage méthodologique. Aucun contenu de l'évaluation n'est annoncé.}

\begin{methodbox}
\textbf{Cinq règles, et rien d'autre.} Ce cadrage ne porte sur aucun contenu de l'évaluation.
\begin{enumerate}
\item \textbf{Gérer le temps.} Trois parties : environ 10 min, 20 min et 15 min. Un regard sur l'horloge à la fin de chaque partie suffit.
\item \textbf{Justifier.} Écrire la relation, la propriété ou la règle avant le calcul. Un résultat seul ne montre pas ce que tu sais.
\item \textbf{Contrôler.} Avant de passer à la suite : ordre de grandeur, substitution, unité, ou vérification par une seconde voie.
\item \textbf{Lire jusqu'au bout.} Souligner la question réellement posée et les données effectivement fournies.
\item \textbf{Passer et revenir.} Une question laissée de côté puis reprise vaut mieux qu'une réponse écrite au hasard.
\end{enumerate}
\end{methodbox}

\begin{graybox}
\textbf{La certitude.} Elle est demandée à trois moments seulement, comme au test de positionnement de début de stage :
\conf\ . Elle n'entre pas dans la note. Elle sert à distinguer ce que tu sais de ce que tu crois savoir.
\end{graybox}

\subsectiontitle{Mon engagement d'avant-évaluation}
\par\vspace{1mm}
\noindent\textbf{Le contrôle que j'écrirai systématiquement :}\ \rule{\dimexpr\linewidth-76mm\relax}{0.32pt}
\par\vspace{3mm}
\noindent\textbf{La question que je traiterai en dernier :}\ \rule{\dimexpr\linewidth-68mm\relax}{0.32pt}
\par\vspace{1mm}

\begin{warningbox}\textbf{Fin du livret de travail.} L'évaluation finale est un document séparé, remis par l'enseignant à l'issue de ces 75 minutes.\end{warningbox}
\end{document}
